\section*{Capítulo \ref{ch:b2:amines} --- Aminas e compostos nitrogenados}

\begin{solution}{exo:b2:amines:1}
Primária; secundária; terciária; sal de amônio quaternário; primária (\omterm{def:b2:huckel:aromatic}{aromática}).
\end{solution}

\begin{solution}{exo:b2:amines:2}
Etanamida $<$ anilina $<$ amônia $<$ metilamina ($\mathrm pK_a$ dos ácidos conjugados
$-0.6$, 4.60, 9.26, 10.66).
\end{solution}

\begin{solution}{exo:b2:amines:3}
\textit{N}-metilciclo-hexanimina, \ce{C6H10=N-CH3}: amina e cetona em meio levemente ácido
(pH 4--5), com remoção da água formada (agente secante ou destilação azeotrópica).
\end{solution}

\begin{solution}{exo:b2:amines:4}
Bromobenzeno; benzonitrila; fenol; 4-hidroxiazobenzeno \ce{C6H5-N=N-C6H4-OH} (acoplamento em para do
grupo \ce{O-} do fenóxido).
\end{solution}

\begin{solution}{exo:b2:amines:5}
$\mathrm pK_a = 9.80$: a pH 7 a razão $[\ce{(CH3)3NH+}]/[\ce{(CH3)3N}] = 10^{2.8} \approx 630$: a
forma protonada predomina; o suco de limão (pH perto de 2) a protona praticamente toda. Agitar uma
solução orgânica de uma amina com ácido aquoso diluído transfere a amina, na forma de seu íon
amônio, para a água; uma base a libera de novo.
\end{solution}

\begin{solution}{exo:b2:amines:6}
1-(ciclo-hex-1-en-1-il)pirrolidina. Após a adição e a perda de água, o nitrogênio leva uma
carga positiva (um íon imínio) e nenhum hidrogênio a perder; em vez disso, um próton é perdido
pelo carbono vizinho, formando a ligação \ce{C=C}.
\end{solution}

\begin{solution}{exo:b2:amines:7}
Benzaldeído com propan-2-amina, ou propanona com benzilamina, cada vez com \ce{NaBH3CN} a pH
perto de 6.
\end{solution}

\begin{solution}{exo:b2:amines:8}
A propiofenona (1-fenilpropan-1-ona), \ce{C6H5COCH2CH3}: uma única adição dá o ânion \omterm{def:b2:amines:imine}{imina},
hidrolisado à cetona.
\end{solution}

\begin{solution}{exo:b2:amines:9}
Ftalimida de potássio + 1-bromo-hexano: \textit{N}-hexilftalimida; depois a hidrazina (ou uma
hidrólise) libera a hexan-1-amina. O nitrogênio da ftalimida alquilada não tem par não ligante
disponível (ele está deslocalizado em duas carbonilas) nem mais \ce{N-H}: ele não pode ser
alquilado uma segunda vez.
\end{solution}

\begin{solution}{exo:b2:amines:10}
A água de bromo broma a anilina nas três posições orto e para do grupo
amino: 2,4,6-tribromoanilina. A \omterm{def:b2:amines:diazonium}{diazotação}, seguida de \ce{H3PO2}, substitui o grupo amino por H: os três
bromos ficam então em 1,3,5 uns em relação aos outros.
\end{solution}

\begin{solution}{exo:b2:amines:11}
4-nitroanilina + \ce{NaNO2} + \ce{HCl} a 0--\qty{5}{\degreeCelsius}: cloreto de
4-nitrobenzenodiazônio. Acoplamento com o fenol em solução fracamente básica (fenóxido, fortemente
ativado) na posição para de \ce{O-}: 4-[(4-nitrofenil)diazenil]fenol, \ce{O2N-C6H4-N=N-C6H4-OH}, um
corante laranja.
\end{solution}

\begin{solution}{exo:b2:amines:12}
O grupo trimetilamônio é um grupo de saída ruim e volumoso, e torna os hidrogênios $\beta$ do
lado \ce{CH3} mais ácidos e mais acessíveis; a base retira um hidrogênio do carbono menos
substituído: but-1-eno, o alceno menos substituído (produto de Hofmann).
\end{solution}

\begin{solution}{pb:b2:amines:1}
\textbf{1.} $\mathrm{{}^-O_3S{-}C_6H_4{-}NH_3^+}$: o grupo sulfônico, ácido, protonou a amina.
\textbf{2.} O zwitteríon é pouco solúvel em água; o carbonato desprotona o grupo amônio,
dando o sulfanilato de sódio, solúvel.
\textbf{3.} \ce{NaNO2 + HCl -> HNO2 + NaCl}.
\textbf{4.}
\[
  \mathrm{{}^-O_3S{-}C_6H_4{-}NH_2 + HNO_2 + H^+ \to {}^-O_3S{-}C_6H_4{-}N_2^+ + 2\,H_2O} .
\]
\textbf{5.} O sal de diazônio se decompõe (em fenol e \ce{N2}) a quente; o ácido nitroso
também se decompõe.
\textbf{6.} $5.00/173.19 = \qty{28.9}{mmol}$: \qty{28.9}{mmol} de \ce{NaNO2}, \qty{1.99}{g}.
\textbf{7.} O ácido nitroso em excesso reagiria com o parceiro de acoplamento; ele é detectado com papel
amido-iodeto e destruído com ureia ou ácido sulfâmico.
\textbf{8.} Seu anel é fortemente ativado pelo grupo dimetilamino, um doador poderoso.
\textbf{9.} Em para de \ce{N(CH3)2}: orientação orto/para, e as posições orto são impedidas pelos
grupos metila.
\textbf{10.} $\mathrm{{}^-O_3S{-}C_6H_4{-}N{=}N{-}C_6H_4{-}N(CH_3)_2}$.
\textbf{11.} Em ácido forte o grupo dimetilamino é protonado: \ce{-NH(CH3)2+} desativa o anel
e o acoplamento para.
\textbf{12.} O produto se forma na sua forma ácida protonada, parcialmente solúvel; em meio básico o
sulfonato de sódio, menos solúvel na solução salina, precipita.
\textbf{13.} Sua carga positiva está deslocalizada no anel e o nitrogênio terminal é um centro
eletrofílico fraco; só reagem os anéis ativados por \ce{-O-}, \ce{-OH} ou \ce{-NR2}.
\textbf{14.} Amarelo-alaranjado em meio básico, vermelho em meio ácido.
\textbf{15.} Em um nitrogênio do grupo azo (a forma protonada é estabilizada por conjugação com
o grupo dimetilamino).
\textbf{16.} Cerca de $\mathrm pK_a \pm 1$: de 3.1 a 4.4 aproximadamente, em torno de $\mathrm pK_a = 3.47$.
\textbf{17.} Dois anéis e o grupo azo formam um longo sistema conjugado: a diferença entre o nível
$\pi$ ocupado mais alto e o vazio mais baixo é pequena (\cref{prop:b2:huckel:gap}) e a absorção
fica no visível.
\textbf{18.} A protonação muda a distribuição de carga no sistema $\pi$ e estreita ainda mais a
diferença: a absorção se desloca para comprimentos de onda maiores, a cor do amarelo ao vermelho.
\textbf{19.} \qty{173.19}{g/mol}.
\textbf{20.} \qty{327.33}{g/mol}.
\textbf{21.} \qty{28.87}{mmol}.
\textbf{22.} $0.02887 \times 327.33 = \qty{9.45}{g}$.
\textbf{23.} Decomposição do sal de diazônio, acoplamento incompleto, solubilidade do produto na
água-mãe e nas águas de lavagem.
\textbf{24.} $0.80 \times 9.45 = \boldsymbol{\qty{7.56}{g}}$ de alaranjado de metila.
\end{solution}
